\section{Data}
\label{sec:data}

% --- fields ---
Our primary field is ECDFS (Rubin tract 5063), where we assemble
$\sim$790 paired Rubin$+$\textit{Euclid} tiles. \todo{exact tile count,
sky area, depth, epoch.} For the out-of-distribution test
(\S\ref{sec:ood}) we use 72 paired tiles in the EDF-S field
(tract 2394).

% --- bands and resolution ---
Each tile carries six Rubin bands ($ugrizy$) and four \textit{Euclid} bands
(VIS, NISP $Y$, $J$, $H$). \textit{Euclid} imagery is taken from the MER
mosaics; VIS and NISP are sampled at $0.1\arcsecpx$ ($\sim$1084$\times$1084
per tile), and Rubin at \todo{Rubin pixel scale / tile shape (6$\times$512$\times$512)}.
The native pixel scales and per-band World Coordinate Systems are retained;
the foundation model resamples internally to a common fused scale of
$0.4\arcsecpx$ (\S\ref{sec:methods}).

\begin{figure*}
  \centering
  \includegraphics[width=\linewidth]{random_10bandtile.png}
  \caption{The ten input channels for a single matched tile. Top: Rubin
    $ugrizy$ ($512^2$ at $0.2\arcsecpx$). Bottom: \textit{Euclid} VIS and NISP
    $Y\!J\!H$ ($\sim$1084$^2$ at $0.1\arcsecpx$, from the MER mosaics). The
    sharp VIS channel and the depth and noise differences across bands are the
    raw material the foundation model fuses.}
  \label{fig:tile}
\end{figure*}

% --- variance / masks ---
Per-pixel variance maps accompany the Rubin imagery and, in the primary field,
the \textit{Euclid} imagery. \todo{Note: the EDF-S \textit{Euclid} tiles used
for the first OOD pass lacked variance maps; see \S\ref{sec:ood}. State final
data state after real-variance tiles are obtained.}

% --- truth / external catalogs ---
For astrometric validation we use \textit{Gaia} \todo{DR} as an absolute
reference, and \todo{the \textit{Euclid} MER catalog for detection
completeness/purity}. We additionally construct a controlled
source-injection set (\S\ref{sec:results}) to measure positional error against
known truth.
